\chapter{Serverumgebung}
\label{app:serverumgebung}

Dieser Anhang dokumentiert die für die Experimente in Kapitel~\ref{cha:realisierung} verwendete Serverumgebung in der Form, wie sie zum Zeitpunkt der Inbetriebnahme bereitgestellt und im Verlauf der Arbeit eingerichtet wurde.
Ziel ist die nachvollziehbare Beschreibung der konkreten Hardware- und Software-Konfiguration, sodass die in Kapitel~\ref{cha:diskussion} dargestellten Resultate auf einer vergleichbaren Maschine reproduziert werden können.

\section{Bereitstellung}
\label{app:serverumgebung_bereitstellung}

Die virtuelle Maschine wurde von HSLU Lab Services im Rahmen der Bachelorarbeit für den Autor bereitgestellt.
Die wesentlichen Eckdaten der Bereitstellung sind in Tabelle~\ref{tab:serverumgebung_provisionierung} zusammengefasst.

\begin{table}[H]
\centering
\caption{Eckdaten der bereitgestellten virtuellen Maschine.}
\label{tab:serverumgebung_provisionierung}
\begin{tabularx}{\textwidth}{l>{\raggedright\arraybackslash}X}
\hline
Attribut & Wert \\
\hline
Bereitsteller & HSLU Lab Services \\
Hostname & \texttt{shadowsrv-001.prod.projects.ls.eee.intern} \\
Zugriff & \gls{ssh}, ausschliesslich aus dem HSLU-Netzwerk (Campus-WLAN oder VPN) \\
Ausserbetriebnahme & 1.~August~2026 \\
\hline
\end{tabularx}
\end{table}

Beim ersten Login auf der VM wurde das initiale Passwort des Administratorkontos in ein nur dem Autor bekanntes Passwort geändert.
Die initialen Zugangsdaten werden in dieser Arbeit aus naheliegenden Sicherheitsgründen nicht wiedergegeben.

Die VM teilt sich gemäss Auskunft von Lab Services ein Subnetz mit weiteren projektgebundenen virtuellen Maschinen anderer Studierender.
Aus diesem Grund wurde der \gls{ssh}-Dienst auf Public-Key-Authentifizierung umgestellt und das Passwort-Login nach dem ersten Login deaktiviert.
Über die VM hinaus sind keine weiteren Dienste nach aussen exponiert.

\section{Hardware}
\label{app:serverumgebung_hardware}

Tabelle~\ref{tab:serverumgebung_hardware} zeigt die zugesagten Hardware-Ressourcen.
Diese Werte gelten für die VM in ihrer Gesamtheit, also inklusive Betriebssystem und Hilfsdienste.

\begin{table}[H]
\centering
\caption{Hardware-Ausstattung der virtuellen Maschine.}
\label{tab:serverumgebung_hardware}
\begin{tabular}{ll}
\hline
Komponente & Wert \\
\hline
vCPU & 32~Kerne \\
Arbeitsspeicher & 128~GB \\
Speicher & 156~GB (LVM, \texttt{ubuntu-vg/ubuntu-lv}) \\
\hline
\end{tabular}
\end{table}

Die Belegung des Root-Dateisystems mit zunächst 33~GB beinhaltet das Basis-Ubuntu-System sowie die in Abschnitt~\ref{app:serverumgebung_software} aufgelisteten Software-Komponenten.
Während eines Simulationslaufes wächst die Belegung in erster Linie durch die \gls{pcap}-Aufzeichnungen.
Bei voller Aufzeichnung aller Eingangs- und Ausgangsknoten wären die 123~GB freier Speicher schnell erschöpft.

\section{Software}
\label{app:serverumgebung_software}

Auf der VM ist das Betriebssystem Ubuntu in einer aktuellen Long-Term-Support-Variante installiert.
Tabelle~\ref{tab:serverumgebung_software} listet die für die Experimente relevanten Software-Komponenten und ihre Quellen.

\begin{table}[H]
\centering
\caption{Software-Komponenten auf der virtuellen Maschine.}
\label{tab:serverumgebung_software}
\begin{tabularx}{\textwidth}{l>{\raggedright\arraybackslash}X>{\raggedright\arraybackslash}X}
\hline
Komponente & Quelle & Verwendung \\
\hline
Ubuntu Server & Bereitstellung Lab Services & Basisbetriebssystem \\
Shadow & GitHub \texttt{shadow/shadow} & Netzwerksimulator \\
Tor & Standardpaket / Tor-Projekt & Tor-Daemon in Shadow \\
Tornettools & GitHub \texttt{shadow/tornettools} & Generierung der Shadow-Konfiguration \\
\texttt{wget2} (gepatcht) & GitLab + Patch aus \texttt{explainwf-popets2023} & Webclient mit \gls{socks}-Unterstützung \\
\texttt{libzim} (Python) & PyPI & Lesen der ZIM-Datei \\
\texttt{zimsrv.py} & Eigene Entwicklung & Webserver in der Simulation \\
Python (Toolsenv) & Bereitstellung über \texttt{toolsenv} & Laufzeit für die Hilfsskripte \\
\hline
\end{tabularx}
\end{table}

Die einmalige Einrichtung der genannten Komponenten erfolgt über das Skript \texttt{setup-wf-server.sh}, das die Abhängigkeiten installiert, \texttt{wget2} mit dem \gls{socks}-Patch baut und die Hilfsskripte (\texttt{newnym.py}, \texttt{zimsrv.py}, \texttt{generate-urls.py}) auf die VM kopiert.

\section{Datenablage und Verzeichnisstruktur}
\label{app:serverumgebung_struktur}

Innerhalb des Heimverzeichnisses des Administratorkontos wird die in Tabelle~\ref{tab:serverumgebung_struktur} dargestellte Verzeichnisstruktur verwendet.
Die genauen Pfade sind in den Hilfsskripten als Konstanten hinterlegt und können bei Bedarf angepasst werden.

\begin{table}[H]
\centering
\caption{Verzeichnisstruktur auf der VM.}
\label{tab:serverumgebung_struktur}
\begin{tabularx}{\textwidth}{l>{\raggedright\arraybackslash}X}
\hline
Pfad & Inhalt \\
\hline
\texttt{\textasciitilde/tornettools} & Installation und Arbeitsverzeichnis von Tornettools \\
\texttt{\textasciitilde/toolsenv} & Python-Virtualenv für Tornettools, libzim und Hilfsskripte \\
\texttt{\textasciitilde/wikidata} & ZIM-Datei der Simple-English-Wikipedia \\
\texttt{\textasciitilde/wget2\_noinstall} & Mit \gls{socks}-Patch gebautes \texttt{wget2} \\
\texttt{\textasciitilde/exp-*} & Pro Experiment ein Unterverzeichnis mit \texttt{shadow.config.yaml}, \texttt{schedule.json}, Logs und Pcaps \\
\hline
\end{tabularx}
\end{table}

Beim Pull der Resultate auf die Arbeitsstation des Autors wird pro Experiment dasselbe Unterverzeichnisschema gespiegelt, sodass lokal und auf der VM dieselbe Struktur vorliegt.
Damit ist die manuelle Nachverfolgung von Resultaten zu ihren ursprünglichen Konfigurationen ohne zusätzliche Werkzeuge möglich.

\section{Ressourcenauslastung der Tamaraw-Läufe}
\label{app:serverumgebung_perfmon}

Die in Abschnitt~\ref{sec:eval_performance} gezeigte Abbildung~\ref{fig:eval_perfmon_comparison} fasst den Ressourcenverlauf der drei Tamaraw-Läufe grafisch zusammen.
Tabellen~\ref{tab:app_perfmon_tamaraw20v2} bis~\ref{tab:app_perfmon_openworld} dokumentieren die zugrunde liegenden Stichproben im Zwei-Stunden-Raster sowie den jeweils letzten von \texttt{perfmon.py} geschriebenen Wert.
Die Werte stammen aus den \texttt{CSV}-Dateien \texttt{perfmon\_tamaraw20v2.csv}, \texttt{perfmon\_tamaraw80.csv} und \texttt{perfmon\_openworld.csv}, die im Quellverzeichnis \texttt{src/thesis/data/} der Arbeit abgelegt sind.
Für \texttt{exp-tamaraw-20v2} liegt keine Disk-Telemetrie vor, da \texttt{perfmon.py} während dieses Laufs noch nicht aktiv war.

\begin{table}[H]
\centering
\caption{\texttt{exp-tamaraw-20v2}: CPU- und RAM-Auslastung der virtuellen Maschine zum jeweiligen Zeitpunkt seit Experimentstart.}
\label{tab:app_perfmon_tamaraw20v2}
\begin{tabular}{rrr}
\hline
Zeit [h] & CPU [\%] & RAM [\%] \\
\hline
 0.00 &   3.20 &  0.82 \\
 2.00 &  99.81 & 24.71 \\
 4.00 &  99.91 & 42.53 \\
 6.00 &  99.78 & 59.93 \\
 8.00 &  99.91 & 77.47 \\
 8.12 &  99.91 & 78.42 \\
\hline
\end{tabular}
\end{table}

\begin{table}[H]
\centering
\caption{\texttt{exp-tamaraw-80}: CPU-, RAM- und Disk-Auslastung der virtuellen Maschine zum jeweiligen Zeitpunkt seit Experimentstart.}
\label{tab:app_perfmon_tamaraw80}
\begin{tabular}{rrrr}
\hline
Zeit [h] & CPU [\%] & RAM [\%] & Disk [\%] \\
\hline
  0.00 & 99.89 & 34.92 & 46.76 \\
  2.01 & 99.87 & 44.68 & 55.63 \\
  4.02 & 99.86 & 54.10 & 64.28 \\
  5.97 & 99.87 & 63.13 & 72.74 \\
  7.98 & 99.88 & 72.34 & 81.35 \\
  9.99 & 99.84 & 81.52 & 89.90 \\
 11.77 & 99.85 & 89.76 & 97.35 \\
\hline
\end{tabular}
\end{table}

\begin{table}[H]
\centering
\caption{\texttt{exp-tamaraw-openworld}: CPU-, RAM- und Disk-Auslastung der virtuellen Maschine zum jeweiligen Zeitpunkt seit Experimentstart. Der Lauf wird bei 16.10 Stunden durch \texttt{ENOSPC} abgebrochen, als die Disk die volle Belegung von 156\,GB erreicht.}
\label{tab:app_perfmon_openworld}
\begin{tabular}{rrrr}
\hline
Zeit [h] & CPU [\%] & RAM [\%] & Disk [\%] \\
\hline
  0.00 &  5.68 &  1.81 & 24.73 \\
  2.02 & 99.86 & 24.61 & 34.16 \\
  3.97 & 99.86 & 34.75 & 43.41 \\
  5.99 & 99.86 & 45.34 & 52.93 \\
  8.01 & 99.82 & 55.62 & 62.41 \\
 10.03 & 99.84 & 65.84 & 71.79 \\
 11.98 & 99.81 & 75.72 & 80.74 \\
 14.00 & 99.84 & 85.87 & 90.08 \\
 16.02 & 99.82 & 95.94 & 99.41 \\
 16.10 & 99.80 & 96.41 & 99.79 \\
\hline
\end{tabular}
\end{table}
